\noindent\hspace*{\parindent}Ich beginne mit der (fingierten) Annahme, dass durch den Zusammenschluss einer sehr großen Zahl von Ehepaaren aus den verschiedensten Altersstufen eine Gesellschaft gebildet worden sei. Alljährlich wird eine Anzahl von Ehen durch den Tod des einen oder des anderen Ehepartners aufgelöst. Dieser Abgang werde dadurch ersetzt, dass jährlich eine bestimmte Zahl neuer Ehepaare hinzukommt, und zwar so viele, dass der Bestand der Gesellschaft insgesamt unverändert bleibe. Von den Witwen, die im Lauf eines Jahres durch den Tod des Ehemanns entstehen, wird, da die Gesellschaft als sehr groß vorausgesetzt wird, ein verhältnismäßig sehr kleiner Teil schon während desselben Jahres wieder sterben. Dadurch verändert sich die Anzahl der am Jahresende tatsächlich vorhandenen Witwen etwas. Am Ende des zweiten Jahres werden ungefähr ebenso viele neue Witwen hinzugekommen sein; dagegen wird ein Teil der Witwen aus dem ersten Jahr bereits wieder gestorben sein, sodass die Witwenzahl am Ende des zweiten Jahres nicht ganz doppelt so groß sein wird wie am Ende des ersten. Am Ende des dritten Jahres sind wieder ungefähr ebenso viele neue Witwen hinzugekommen; dagegen wird der zu Beginn des dritten Jahres vorhandene Bestand einen fast doppelt so großen Abgang erlitten haben, wie ihn die Bilanz des zweiten Jahres ergeben hatte. Man sieht, dass die Witwenzahl auf diese Weise zwar fortwährend wächst, aber immer langsamer, bis sie schließlich so groß geworden ist, dass der jährliche Abgang durch den Tod der Witwen dem Zugang durch den Tod von Ehemännern aus der Gesellschaft das Gleichgewicht hält. Dann wird also der Beharrungszustand eintreten: Die Witwenzahl hat ihr Maximum erreicht und bleibt von da an unverändert.

Fügen wir den obigen Voraussetzungen noch hinzu, dass sowohl die Anzahl der in jedem Jahr
\SourceLeaf{146}
neu hinzukommenden Ehepaare als auch das Verhältnis, in dem in einer solchen Gruppe die verschiedenen Altersstufen gemischt sind, Jahr für Jahr gleich bleiben und ebenso die Sterbefälle genau den Sterblichkeitstafeln entsprechen, so wird jener Beharrungszustand seinen Namen in aller Strenge verdienen. Sowohl die Anzahl der Ehepaare in der Gesellschaft wird völlig unverändert bleiben als auch die Anzahl der Witwen, nachdem diese einmal ihren höchsten Wert erreicht hat. Ferner ist klar: Stellt man sich eine zweite ähnliche Gesellschaft vor, in der die neu Beitretenden in genau demselben Verhältnis gemischt sind wie in der ersten, die aber doppelt so viele Ehepaare umfasst, so wird auch die höchste Witwenzahl doppelt so groß sein. Allgemein ausgedrückt wird das Verhältnis der bestehenden Ehen zur Witwenzahl nach Eintritt des Beharrungszustands nicht vom Umfang der Vereinigung abhängen, sofern dieser nur als konstant bleibend betrachtet wird, sondern allein vom Verhältnis, in dem die verschiedenen Altersstufen der neu Beitretenden gemischt sind. Wäre dieses letztere Verhältnis vollständig bekannt, so ließe sich das erstere a priori durch Rechnung bestimmen. Der Beharrungszustand wird übrigens, wenn nicht früher, jedenfalls dann eintreten, wenn die ursprünglichen Teilnehmer alle gestorben sind. Wenn man die äußersten Fälle berücksichtigen will, könnte sich dies möglicherweise bis 80 Jahre nach dem ersten Zusammenschluss verzögern. Man kann jedoch annehmen, dass der wirkliche Zustand bereits nach 45—50 Jahren dem Beharrungszustand sehr nahe gekommen sein wird\footnote{\textsc{Kritter} begnügt sich in seinem Gutachten schon mit 40 Jahren.}. Um die Vorstellungen bestimmter zu machen, will ich beispielshalber konkrete Zahlen nennen. Da diese jedoch nur zur Erläuterung des sonst abstrakten Vortrags dienen sollen, erheben sie keinen Anspruch auf vollkommen genaue Angemessenheit. Die Gesellschaft bestehe aus 2600 Ehepaaren, zu denen jährlich 130 neue hinzukommen, während durchschnittlich ebenso viele ausscheiden, und zwar 70 durch den Tod des Mannes, 60 durch den Tod der Frau. Von den 70 Witwen, die im Lauf eines Jahres entstehen, stirbt eine schon im selben Jahr, sodass am Ende des ersten Jahres 69 Witwen vorhanden sind. Von diesen sind am Ende des zweiten Jahres weitere 2 gestorben; dagegen sind wieder 69 neue Witwen hinzugekommen, sodass insgesamt 136 vorhanden sind. Von diesen sterben im dritten Jahr 4, und da wieder 69 neue hinzukommen, beträgt die Gesamtzahl am Ende des dritten Jahres 201. Auf diese Weise immer langsamer fortschreitend mag die Witwenzahl nach 10 Jahren 600, nach 20 Jahren 1000, nach 30 Jahren 1200, nach 40 Jahren 1280, nach 45 Jahren 1294 und nach 50 Jahren 1298 betragen; später kommen noch ein paar hinzu und ergeben das wirkliche Maximum von 1300. Hier hätten wir also das Verhältnis der bestehenden Ehen zur höchsten Witwenzahl wie zwei zu eins.

In einer wirklichen Gesellschaft, in der die geforderten Bedingungen nicht in ihrer genauen Strenge, sondern nur durchschnittlich gelten, kann es natürlich nicht so regelmäßig zugehen wie in der fingierten. In jener werden jährlich nicht genau 130 neue Ehepaare beitreten, sondern in einem Jahr etwas mehr, in einem anderen etwas weniger. Ebenso werden die verschiedenen Altersstufen der neu Beitretenden in einem Jahr etwas anders gemischt sein als in einem anderen oder als in der idealen Gesellschaft angenommen. So wird etwa das Durchschnittsalter der Männer oder das der Frauen oder der durchschnittliche Unterschied zwischen beiden einmal etwas größer, ein anderes Mal etwas kleiner sein. Schließlich werden auch bei einer gegebenen Personenzahl die Sterbefälle nicht \textit{genau} den Sterblichkeitstafeln folgen. Vielmehr werden diese in einem Jahr etwas zu wenig, in einem anderen etwas zu viel angeben. All dies wird dazu führen, dass die Witwenzahl in der wirklichen Gesellschaft zu einer Zeit etwas größer sein wird als in der idealen, zu einer anderen etwas kleiner. Ein eigentlicher Beharrungszustand \textit{im strengsten Sinn} wird in jener niemals eintreten, sondern ein Schwanken um einen Mittel%
\SourceLeaf{147}%
zustand. Zu der Zeit also, zu der die Witwenzahl in der idealen Gesellschaft das Maximum 1300 ganz oder fast ganz erreicht hätte, wird es in der wirklichen ein Auf- und Abschwanken über und unter diese Zahl geben. Zu einer Zeit können 1350, wohl auch 1360 vorhanden sein, zu einer anderen 1250, wohl auch nur 1240. Es ist aber ganz unmöglich, hier genaue Grenzen zu setzen und von irgendeiner Zahl, welche man auch wähle, mit Bestimmtheit zu behaupten, diese könne zwar noch erreicht werden, die nächsthöhere aber nicht mehr. Die Wahrscheinlichkeitsrechnung lehrt in ihrem heutigen Zustand, solche Schwankungen auf ein bestimmtes Maß zurückzuführen. Dieses wird jedoch nicht vom Extrem hergenommen, sondern aus der Berücksichtigung aller Zwischenstufen und ihrer relativen Wahrscheinlichkeit. Natürlich ist hier keineswegs der Ort, diese Theorie weiterzuverfolgen. Ihre Anwendbarkeit auf konkrete Fälle hängt außerdem davon ab, dass man entweder eine mathematisch genaue Kenntnis der bedingenden Elemente hat, wie es etwa bei Glücksspielen gewöhnlich der Fall ist, oder dass genügend reichhaltige Erfahrungen verfügbar sind. Bei Gegenständen wie Witwenkassen trifft beides nicht zu, sodass a priori für eine solche Bemessung nichts getan werden kann. Die Theorie lehrt aber doch, aus den Schwankungen, die in einem Fall vorgekommen sind, auf diejenigen zu schließen, die man mit gleichem Recht in einem qualitativ gleichen, aber quantitativ verschiedenen anderen Fall erwarten muss. Ihrer absoluten Größe nach werden die Schwankungen in einer großen Gesellschaft größer sein als in einer kleinen; bei ihrer relativen Größe wird es sich umgekehrt verhalten. Hätte also etwa eine Gesellschaft der obigen Art zu einer Zeit bis zu 60 Witwen mehr, zu einer anderen bis zu 60 weniger gezählt als der Durchschnitt von 1300, und wäre dabei die Zahl 1360 oder eine ihr nahekommende mit einiger Andauer vorgekommen und ebenso die Zahl 1240 oder eine ihr nahe liegende, so dürfte man schließen, dass in einer anderen Gesellschaft, die bei sonst ähnlichen Verhältnissen durchschnittlich nur 26 Ehepaare zählt, die Normalwitwenzahl ebenso leicht um 6 vermehrt oder vermindert erscheinen könne. Etwas anders ausgedrückt: Ebenso gut wie die Witwenzahl in der großen Gesellschaft zwischen 1360 und 1240 kann sie in der kleinen zwischen 7 und 19 auf- und abschwanken. Das \textit{absolute} Schwanken ist, wenn der Umfang der großen Gesellschaft 100 Mal größer ist als der Umfang der kleinen, in der großen zehnmal so groß wie in der kleinen. Beim \textit{relativen} Schwanken, das in der großen Gesellschaft nach den obigen Zahlen \(4\tfrac8{13}\) Prozent betragen würde und in der kleinen \(46\tfrac2{13}\) Prozent, ist das Verhältnis genau umgekehrt.

\Needspace{24\baselineskip}
So viel zur Sache. Was die Bezeichnung betrifft, habe ich eben für die Zahlen der Beispiele 1300 und 13 statt der weitschweifigen Umschreibung ‘Durchschnittszahl der Witwen nach Eintritt des Beharrungszustands’ die Bezeichnung ‘Normalwitwenzahl’ gebraucht, die man leicht als nicht unpassend gelten lassen wird. Meines Wissens ist sie aber nicht \textit{üblich}. Es ist vielmehr ganz gewöhnlich, jenen Begriff kurz als \textit{höchste Witwenzahl} zu bezeichnen. Dabei sieht man von den unregelmäßigen Schwankungen völlig ab und berücksichtigt nur den Beharrungszustand im Allgemeinen im Gegensatz zur vorhergegangenen Zeit. Für jemanden, der einigermaßen sachkundig ist, wird ein Missverständnis hierbei nicht leicht möglich sein. Jedenfalls ist aber so viel gewiss und aus dem Gutachten, wenn man es nur mit einiger Aufmerksamkeit liest, sofort zu erkennen, dass \textsc{Kritter} darin die höchste Witwenzahl in \textit{diesem} und nur in diesem Sinn verstanden hat. P. hingegen dachte sich dabei etwas ganz anderes, nämlich die Zahl, über die die Anzahl der Witwen nach höchster Wahrscheinlichkeit niemals hinausgehen können sollte. Beide Zahlen zu vermengen ist ungefähr dasselbe, als wenn man den Durchschnittspreis einer Handelsware mit ihrem höchsten Preis verwechselte. Damit ist der oben S. [140] unter 3 gemachte Einwand hinreichend bewiesen. Eine Frage nach der höchsten Witwenzahl in P.s Sinn des Wortes ist jedoch eine solche, auf die sich überhaupt keine bestimmte Antwort geben lässt.
